\documentclass{article}
\usepackage{amsmath,url}
\title{What planned bridge activation guarantees before an abrupt stop}
\begin{document}
\maketitle
\paragraph{Matched operating regimes.}
The native baseline is \emph{Safe Formation Control of Open Multi-Robot Systems with Connectivity-Preserving Reconfiguration} by \c{S}ekercio\u{g}lu, De Carli and Dimarogonas, \url{https://arxiv.org/html/2609.29939v1}. Its manager maintains the departing robot's old links until the prospective bridge is ready, then removes the old links. The paper-level double-integrator model assumes that the desired acceleration is realized. A crash at request time instead removes communication immediately and leaves a body whose exclusion radius still applies. The former result is not a claim about the latter. Comparable named literature: Maeda, Endo and Matsuno's decentralized navigation after failure (DOI 10.1016/j.jfranklin.2023.05.013) was verified at the abstract level; its full controller was not available for this replay. Breeden and Panagou's composed barriers require the intersection to be a viability domain (\url{https://arxiv.org/abs/2210.01354}); instantaneous feasible input, graph connectivity, and subsequent trajectory safety are separate questions.

\paragraph{Author-supplied point-mass cases.}
Run the original Python controller and simulation, with the point-mass gains in its example, on both default scenarios supplied by the authors. No outcome was used to select cases: inspect all four scheduled leave requests and all corresponding bridge activations and departures. At every request cut, removing just the departing radio splits even the \emph{full} nominal distance-based graph in two, not just the manager's chosen graph. The smallest distance between components exceeds its nominal radius by \(0.3992,\ 0.1532,\ 0.7036,\ 0.5230\) in order. Thus imposing immediate connectivity in those counterfactual failures is impossible at that cut. At every corresponding managed bridge activation and actual departure, the surviving proposed graph and the full nominal graph are connected. This reproduces the original finite experiment, not the author's crash test and not a held-out mission.

A separate counterfactual takes each departing body at the actual managed departure pose and holds it fixed while replaying the original survivor trajectory. There is no new controller in this calculation. To check between samples, use that the simulator's point-mass step follows constant acceleration. The maximum deviation of this arc from its chord is \(\|a\|\Delta t^2/8\), since the difference is \(a\,t(t-\Delta t)/2\); the norm is 1-Lipschitz in position. None of the checked post-departure active segments is close enough to the frozen body to violate the prescribed clearance under these bounds. This negative stress test is reported explicitly. It neither simulates a crash during the managed bridge nor proves that a controller that knows of the stopped body follows the same trajectory.

\paragraph{Exact speed-only recovery on the request geometries.}
As a different, explicitly simpler abstraction, fix the post-loss positions, stationary body and common speed cap \(V>0\). Assume full pose information, immediate authorized control, no acceleration constraint or disturbance, and two surviving components. Their internal required distances and links must be maintained, collision clearance \(d\) is continuous, and one cross-component bridge must be obtained by a deadline at the closed range threshold \(R_*\!>d\). Set \(s\) to the closest cross-pair distance, \(s>R_*\). Then no policy can bridge before \((s-R_*)/(2V)\), by the length and triangle bounds on motion of the endpoints. Translate each entire component toward the other using directions defined by a closest pair, at speed \(V\). At the lower bound the pair attains \(R_*\); intra-component offsets are unchanged, and all cross distances are at least \(s-2Vt\ge R_*\) over the interval. The only extra obligation is clearance of each straight segment from the failed body. If those segments clear it, the bound is attained; an approach blocked by the body rules out that witness, not every recovery.

For all four original request poses the test passes with \(R_*=\text{nominal range}-\text{source activation margin}\). At unit speed the exact formula gives approximately \(0.2296,\ 0.1166,\ 0.3918,\ 0.3015\) time units, and the straight-segment body-clearance checks have strict slack. These are not runnable timings for the original acceleration-based controller and do not include fault recognition, an edge activation decision, low-level actuation, or radio-link certification. A direct shortest-distance bound with the same information proves the same optimum.

\paragraph{Matched infeasible geometry.}
Take two single-robot components at \((-2,0),(2,0)\), with a failed body at \((-1,0)\), clearance \(1/2\), bridge threshold 2, and unit speed. At time \(1+\epsilon\) the left endpoint must have \(x\in[-1-\epsilon,-1+\epsilon]\) to be within range of a reachable right endpoint. Its travel bound also implies \(y^2\le 4\epsilon\) for \(0\le\epsilon\le1\). Its squared distance from the failed body is at most \(\epsilon^2+4\epsilon\), which is below \(1/4\) at \(1.05\). Thus neither steering nor a more informed controller can meet that deadline, even though bridging without body clearance or remaining safe without bridging is separately feasible. Moving both upward one unit and then together toward the midpoint horizontally one unit gives a checked safe bridge by time 2. No exact optimal detour time is claimed. The scoped result distinguishes necessary immediate failure, continuous safety and deadline restoration. The gap to R6 completion includes acceleration limits and local evidence in the same matched post-crash family, plus the full named crash comparator.
\end{document}
